On dispose d’un sac et de deux urnes A et B.

\begin{itemize}
\item Le sac contient 4 boules : 1 boule avec la lettre A et 3 boules avec la lettre B.
\item L’urne A contient 5 billets : 3 billets de $50$~euros et 2 billets de $10$~euros.
\item L’urne B contient 4 billets : 1 billet de $50$~euros et 3 billets de $10$~euros.
\end{itemize}

\noindent
Un joueur prend au hasard une boule dans le sac :
\begin{itemize}
\item si c’est une boule avec la lettre A, il prend au hasard un billet dans l’urne A.
\item si c’est une boule avec la lettre B, il prend au hasard un billet dans l’urne B.
\end{itemize}

\noindent

On note les évènements suivants :

\begin{itemize}[label=$\bullet~$]
\item $A$ : le joueur obtient une boule avec la lettre A.
\item $C$ : le joueur obtient un billet de 50 euros.
\end{itemize}

\medskip

\begin{minipage}{0.6\linewidth}
\begin{enumerate}
\item Recopier et compléter l'arbre ci-contre représentant la situation.
\item Quelle est la probabilité de l'évènement \og \emph{le joueur obtient une boule avec la lettre {\rm A} et un billet de $50$ }\euro \fg{} ?
\item Démontrer que la probabilité $P(C)$ est égale à \np{0,3375}.
\end{enumerate}
\end{minipage}\hfill
\begin{minipage}{0.38\linewidth}
\begin{center}
\pstree[treemode=R,levelsep=2cm]{\TR{}}
{\pstree{\TR{$~A~$}\taput{\ldots}}
	{\TR{~$C$} \taput{\ldots}
	\TR{~$\overline{C}$} \tbput{\ldots}
	}
\pstree{\TR{~$\overline{A}~$}\tbput{\ldots}}
	{\TR{~$C$} \taput{\ldots}
	\TR{~$\overline{C}$} \tbput{\ldots}
	}
}
\end{center}
\end{minipage}

\begin{enumerate}[start=4]
\item Le joueur a obtenu un billet de 10 euros.

L'affirmation \og \emph{Il y a plus de $80\,\%$ de chances qu'il ait au préalable obtenu une boule avec la lettre} B \fg{} est-elle vraie ? Justifier.
\item On note $X_1$ la variable aléatoire qui donne la somme, en euros, obtenue par le joueur.

Exemple : si le joueur obtient un billet de $50$~\euro, on a $X_1 = 50$.

Montrer que l'espérance $E(X_1)$ est égale à $23,50$ et que la variance $V(X_1)$ est égale à $357,75$.
\item Après avoir remis la boule dans le sac et le billet dans l'urne où il a été pris, le joueur joue une deuxième partie. On note $X_2 $ la variable aléatoire qui donne la somme obtenue par le joueur lors de cette deuxième partie.

On note $Y$ la variable aléatoire ainsi définie : $Y = X_1 + X_2$.
	\begin{enumerate}
		\item Montrer que $E(Y) = 47$.
		\item Expliquer pourquoi on a $V(Y) = V(X_1) + V(X_2)$
	\end{enumerate}
\item Le joueur joue de même une troisième, une quatrième,\ldots, une centième partie.

On définit donc de la même façon les variables aléatoires $X_3, \:X_4,\:\ldots, X_{100}$.

On note $Z$ la variable aléatoire définie par $Z = X_1 + X_2 + \ldots + X_{100}$.

Démontrer que la probabilité que $Z$ appartienne à l'intervalle $]\np{1950}~;~\np{2750}[$ 
 est supérieure ou égale à $0,75$.
\end{enumerate}

\noindent

\bigskip
